\documentclass[../Main.tex]{subfiles}
\begin{document}

\section{What are Alternative Games?}
\label{chap:5}
To develop a concept for an alternative game (as described in the research goal), a clear understanding of which games can be categorized as alternative games are crucial. Let us first take a look at what the word “alternative” means in this context with some definitions from three well-known dictionaries:

\blockquote{
    \textit{“Different from the usual or traditional way in which something is done”} \\
    - Oxford learner's dictionary.
    
    \textit{“Different from the usual or conventional”} \\
    - Merriam-Webster's Collegiate Dictionary.
     
    \textit{“Alternative things are considered to be unusual and often have a small but enthusiastic group of people who support them”} \\
    - Cambridge dictionary.
}

While all of these dictionaries define the word differently, we can see that there are some commonalities. All three describe alternative as something “unusual” or “different from the usual.” The key takeaway from the definitions is that an alternative game is a game where the creators have done something in such a way that it differs from other games.

To recognize alternative game concepts more easily, we have looked at different existing alternative games and tried to find out what they have in common. The result of this research is a list of six characteristics, where alternative games often differ from ordinary ones.

\subsection{Alternative game characteristics}
Our proposal of which games can be considered alternative games is then the following:

\textit{If a game deviates from regular games in one or more of the game characteristics categories mentioned in this chapter, it can be considered an alternative game.}

Below are explanations of all categories with coherent, relevant alternative games presented. These categories are covering all alternative games we have studied during this phase of the project. However, it is important to point out that there might exist additional characteristics from games that we are not aware of or that are not developed yet. As technology moves forward, it is reasonable to assume that new branches of alternative games might emerge.

\subsubsection{Input}
A common way of creating an alternative game is to use an unusual input source (such as the GPS, microphone, special controllers, or accelerometer) for controlling the game. Three examples of games using alternative input sources are Knowledge War \citep{wang_knowledge_2016}, CityZombie \citep{zhu_experiences_2010} and Exermon \citep{wang_evaluation_2017}. Both Knowledge War and CityZombie use the user’s position as a central part of playing the game. Moving around in the real world translates to the gameplay. Exermon is an example of utilizing the accelerometer on the player's phone. The game is a strength-based exercise game that tries to motivate people to do more strength exercises. An alternative input source, the accelerometer, is used to register the number of repetitions the user completes in one workout.

\subsubsection{Output}
The output from the game is another characteristic commonly explored in the development of alternative games. Some known examples are sound, haptics, and augmented reality, but there are many more. Two examples of games using alternative outputs are “Unexceptional.net” and “The Amazing City Game.” Unexceptional.net is a story-driven pervasive game where “the story evolves through multiple media channels interacting with the player.” The game includes different untraditional outputs, such as SMS and phone calls, to provide feedback to the player \citep{wang_unexceptionalnet_2009}. The Amazing City Game is also a pervasive game where the players can play a knowledge competition tour. The game utilizes location and augmented reality to show different vital landmarks at a user’s location \citep{wang_pervasive_2011}.

\subsubsection{Target group}
The target group is an essential characteristic of a game, and it often decides the type of users appealed to the game should. Some games include adjustments to target specific groups of people, and some of these can make the game fall into the alternative game definition. More special target groups, like people with physical disabilities, often require more adjustments to play games at all. For example, people who cannot maneuver a game using traditional handheld controllers require other ways to interact. A Brain-Computer Interface (BCI) could be a way of bringing games to this group of people. One concrete example is the use of brain waves (via a \href{www.neurosky.com}{NeuroSky} brain set) to control the movement in a Snake game \citep{wang_using_2012}.

\subsubsection{Game genre}
Computer games are generally categorized as one of eight game genres \citep{prensky_fun_2001}:

\begin{itemize}
    \item Action Games
    \item Adventure Games
    \item Fighting Games
    \item Puzzle Games
    \item Role-Playing Games (RPG)
    \item Simulation Games
    \item Sports Games
    \item Strategy Games
\end{itemize}

Earlier computer games were usually limited to one category, while today, games often fall into multiple ones. Additional categories may also be defined for games that do not directly belong to any of these genres. Alternative games often consist of unique concepts and uncommon gameplay, often falling into other categories than those mentioned above. An example of such a genre is gamification, defined as “the application of game design principles in non-gaming contexts” \cite{robson_is_2015}. An example of a game not fitting directly into any of the eight general game genre categories, but that utilizes gamification, is the “Kinect Recycling Game.” This game is a gamification concept of recycling trash, where the goal is to teach recycling in an engaging and realistic way. As the name implies, the game uses Kinect to recognize players' hand gestures to place the trash into the right bins \citep{wang_learning_2015}.

\subsubsection{Domain}
Computer games can be divided into different domains. The domain of a game depends on what area of use it serves. Some examples of game domains are entertainment, education, industry, government, and health. Alternative games may sometimes define their area of use, where the new domain may be a combination of other domains or a subset of an existing domain. For instance, edutainment games are a crossing between education and entertainment. Another example of an alternative domain is Attensi’s system for gamified work training, mixing industry, and learning \citep{attensi_why_2020}. The company offers tailored 3D job simulators to large organizations to achieve real and measurable changes in employees’ behavior and skills.

\subsubsection{Interaction}
There exists a lot of various interaction forms for games. Examples included cooperation, competition, player vs. player (PVP), and more. Interaction can also contain an alternative twist, i.e., interaction forms that are not used in ordinary games. One such concept is the multimodal “Multiplayer On One Screen Entertainment System (MOOSES)” approach. Here, the players use their smartphones as controllers for gaming on a shared screen \citep{wang_evaluation_2010}. \href{http://www.airconsole.com}{AirConsole} is a concrete example of where this concept is utilized. The console includes multiple MOOSES based games where the players use their smartphones as controllers. Games can also be pervasive, meaning the game experience is extended beyond the fictive world. The player interacts with both the virtual and real worlds.

\subsection{Framework [Not written yet]}

\subsection{Summary}
This chapter presented our definition of alternative games together with the six coherent categories of game characteristics. A short description of each of them with relevant examples from the literature was given. In the next chapter, existing popular and known alternative games will be described. How they fit into the proposed categories in this chapter will also be outlined.

\biblio % Needed for referencing to working when compiling individual subfiles - Do not remove
\end{document}